\subsection{Source controls, compatibility and complete-operation costs}

The final native audit uses the nineteen-case maintained corpus, the empty
circle and the valid one-component closures retained from 240 seeded random
braids. These give 84 source diagrams. Each is tested in its original form,
as a mirror, and after reversing the crossing list and rotating each PD row
by two slots: 252 diagram variants, with at most 141 crossings.

Every mixed crossing triangulation is isomorphic to Regina's unsimplified
diagram complement and has vertex-link Euler characteristics $2,2,0$.
For each variant, both finite subdivisions pass the native manifold validator
and Regina's finite orientable one-torus-boundary checks: 504 comparisons.
The current source checker accepts both outputs, and the original checker
accepts the retained centred output: 756 successful source replays. All 252
centred outputs exactly equal the first published constructor at commit
\code{8354e8aab}. The external engine is Regina 7.4, distribution 7.4.1.
Neither engine's manifold checks alone would establish knot provenance.

The focused tests disable the producer during independent replay, check
nondegeneracy and exact volume for both local subdivisions, reject all 400
single-face mutations of the two one-crossing outputs, and reject a
reciprocal but incorrect gluing. They also reject a trefoil exterior offered
for a three-crossing unknot source, malformed schemas and unchecked invalid
PDs. Cancellation propagates without input mutation. A source-bound boundary
vertex-link disc correctly has zero essential-disc count; changing that
certificate's count to one is rejected. Seven small external solid-torus
controls cover the empty circle, curls of both signs, a reducible
three-crossing unknot, a trefoil, a figure-eight and a stabilized circle.
The original implementation passed 1,149 maintained tests. The final version
adds the retained-subdivision contract test and passes all 1,150 maintained
tests in 216.795 seconds. Its eight focused tests pass in 1.353 seconds.
The final geometry audit took 56.440 seconds while the suite also ran;
that elapsed time is not used as a performance comparison.

There was no newly discovered normal meridian at this construction-only
checkpoint; Section~\ref{sec:cocycle-seeds} subsequently supplies native
candidates and source-bound disc proofs. Two
exploratory external calls to \code{nonTrivialSphereOrDisc()}, one on the
80-tetrahedron circle and one on the 20-tetrahedron circle, each reached its
20-second process cap without returning a surface. These are incomplete
searches, not negative results or a speed comparison. The existing supplied
normal-vector algorithms can now be applied to source-authenticated geometry,
but obtaining useful vectors efficiently remains an independent obligation.

\subsubsection{Measured effect of the pulling subdivision}

The final benchmark runs serially after the suite, audit and publication of
the runtime change. It compares the pinned first constructor and checker at
\code{8354e8aab} against the final pulling implementation, using the same
maintained manifold validator. Each timed call starts afresh, constructs the
exterior, independently replays its source geometry and validates the finite
manifold. Serialization for the stored output hash is outside the timer.
Both variants perform the complete stated operation; they have different
canonical subdivisions and therefore different expected output hashes.

The four inputs are stabilized circles with $1,8,32,128$ crossings. There
are five rounds and four interleaved arms: old, old repeat, new and new
repeat, plus one warmup per arm and input. All eighty measured calls and
sixteen warmups complete in a total driver elapsed time of 24.675 seconds.
Hashes agree within each subdivision across every repeat. The table gives
unpaired time medians and medians of paired ratios.

\begin{center}
\small
\begin{tabular}{rrrrrrr}
\toprule
Crossings & New tets & Old (ms) & New (ms) & Old/new & Old A/A & New A/A \\
\midrule
1 & 20 & 8.444 & 2.269 & 3.748 & 0.996 & 1.043 \\
8 & 160 & 73.140 & 17.259 & 4.346 & 1.032 & 0.997 \\
32 & 640 & 285.732 & 67.210 & 4.228 & 1.021 & 0.998 \\
128 & 2560 & 1217.252 & 284.316 & 4.299 & 0.998 & 0.996 \\
\bottomrule
\end{tabular}

\end{center}

The complete-call old/new ratios are $3.748$--$4.346$. Old A/A ratios are
$0.996$--$1.032$ and new A/A ratios are $0.996$--$1.043$. At 128 crossings,
the median falls from $1.217252$ seconds to $0.284316$ seconds; tetrahedra
fall from 10,240 to 2,560. This is consistent with the exact factor-four
reduction in geometry. It measures construction, source replay and manifold
validation, not a normal-vector search, a simplified minimal exterior or
whole recognition. It supplies no new general quasipolynomial theorem.


Reproduction uses \path{../fast/normal_orbit_research/exteriors.py}, with
\code{audit} and \code{benchmark} modes. The driver pins all 457 Python and
fixture sources before and after each run, and separately pins the two
baseline modules loaded from Git. Final records have the
\code{diagram-exterior-} prefix in \path{data/}; the first implementation's
audit, suite and A/A timing records retain the
\code{diagram-exterior-centred-} prefix. The coordinate and gluing proofs are
in the maintained article, not delegated to the optional external controls.
